\section{Extracting a small vertex cover}\label{sec:soundness}

\begin{proposition}[Soundness]\label{prop:soundness}
If the instance of Theorem~\ref{thm:reduction} packs into at most
\(B+c\) bins, then its input graph has a vertex cover of size at most
\(k+\rho n\).
\end{proposition}

We first use deadline prefixes to force coverage at every test point.
The competing trees then force a large global allocation on one side
of each vertex. Repeated edge jobs turn the plus allocations into a
vertex cover, whose size is bounded by the global item inventories.

Fix such a packing throughout this section.  By
Lemma~\ref{lem:good-anchors}, at most \(3K\) anchors of any fixed label
fail to belong to good tuples.  This bound applies to every subset of
that label's anchors, including any chosen collection of deadlines or
keys.  Summing over labels gives a bound of \(3nK\) when needed.

Let \(\mathcal M_v\) be the set of good \(\mathrm M\) tuples of label
\(v\).  Each has a row baseline \(b\), a completion coordinate \(h\le b\),
and the associated coverage interval \([h,b)\).
Write \(x_v^+\) and \(x_v^-\) for the numbers of these tuples using,
respectively, a length-one \(\Up\) or \(\Um\) item.  Distinct tuples
use distinct items, so the inventories give
\begin{equation}\label{eq:sound-global-budgets}
 \sum_v x_v^+\le dk,\qquad
 \sum_v x_v^-\le d(n-k),\qquad
 \sum_v(x_v^++x_v^-)\le dn.
\end{equation}
All interval counts below include multiplicities of tuples.

\subsection{Coverage and a forced side}

For a real point \(a\), define
\[
 A_v(a)=\#\{M\in\mathcal M_v:b(M)\le a\},\qquad
 C_v(a)=\#\{M\in\mathcal M_v:h(M)\le a\},
\]
and let \(I_v(a)\) count the intervals of \(\mathcal M_v\) containing
\(a\).  Because \(h\le b\), the identity
\(\ind\{h\le a\}=\ind\{b\le a\}+\ind\{h\le a<b\}\)
holds even at tied endpoints.  Consequently
\begin{equation}\label{eq:sound-prefix-decomposition}
 C_v(a)=A_v(a)+I_v(a).
\end{equation}
The function \(B_v(a)\) remains the original baseline prefix used to
define the deadlines; it need not equal \(A_v(a)\) in this packing.

\begin{lemma}[Coverage at every test point]\label{lem:coverage}
For every label \(v\) and every point \(a\) in one of its half-open
test intervals,
\begin{equation}\label{eq:sound-test-coverage}
 I_v(a)\ge d-6K.
\end{equation}
\end{lemma}

\begin{proof}
There are exactly \(B_v(a)+d\) deadlines at most \(a\).
At least \(B_v(a)+d-3K\) of their anchors are good.
Their distinct \(\mathrm M\) matches have \(h\le z\le a\) by
Lemma~\ref{lem:good-feasibility}, giving
\begin{equation}\label{eq:sound-completion-prefix}
 C_v(a)\ge B_v(a)+d-3K.
\end{equation}
The deadline formula uses the half-open test convention: moving
\(d\) copies from a right endpoint \(r_0\) to a left endpoint
\(\ell_0\) changes the prefix by
\(d(\ind\{\ell_0\le a\}-\ind\{r_0\le a\})
=d\ind\{\ell_0\le a<r_0\}\).

Let \(N'\) be the number of positions in \(\mathcal R_v\) with
\(b_r\le a\).  There are \(2dN'\) actual \(\Srow\) rows at those
positions and \(dN'\) keys there.  At least \(dN'-3K\) of these keys
have good \(\mathrm G\) matches.  The row-order conclusion of
Lemma~\ref{lem:good-feasibility} places all their rows in this same
baseline prefix.  Since these are distinct rows, at most
\[
 2dN'-(dN'-3K)=dN'+3K
\]
of the prefix's \(\Srow\) rows can occur in good \(\mathrm M\)
tuples.  This inequality is valid also when \(dN'-3K<0\).
Other uses of prefix rows only reduce the number available.
The \(\Tp,\Tm\) rows contribute at most their entire original
baseline prefix.  Since \(B_v(a)\) includes exactly \(dN'\)
hypothetical job baselines, we obtain
\begin{equation}\label{eq:sound-baseline-prefix}
 A_v(a)\le B_v(a)+3K.
\end{equation}
Subtracting Equation~\eqref{eq:sound-baseline-prefix} from
Equation~\eqref{eq:sound-completion-prefix} and using
Equation~\eqref{eq:sound-prefix-decomposition} proves the claim.
\end{proof}

\begin{lemma}[A large count on one side]\label{lem:side-choice}
Every vertex \(v\) satisfies
\begin{equation}\label{eq:sound-side-choice}
 x_v^+\ge d-6K\qquad\text{or}\qquad x_v^-\ge d-6K.
\end{equation}
\end{lemma}

\begin{proof}
A length-one global item can contribute coverage only on its own
side.  Indeed, a minus row has baseline at least \(4\), is not short,
and has completion at least \(4-1-\lambda_1>1\).
Every plus row, including each \(\Srow\) row, has baseline at most
\(1\), so its interval cannot cover a minus test point.

Choose one interior point in each tree leaf cell.  A tuple whose
global and local lengths are both zero has an empty interval unless
its row is short.  In the latter case its interval is exactly its
job test interval, which is disjoint from all leaf cells.
For a tuple with global length zero and positive local length
\(\lambda_\ell\), the interval length is at most
\[
 \lambda_\ell+\Delta\le \frac{11}{10}\lambda_\ell<3\lambda_\ell.
\]
By Lemma~\ref{lem:geometry}, it therefore meets at most one closed
depth-\(\ell\) cell across the two trees.  If there is such a cell,
mark its tree vertex.  In particular, if the interval covers a
chosen leaf point, that point belongs to its depth-\(\ell\)
ancestor cell; the marked vertex is this ancestor.
This reasoning depends on the interval's length, regardless of
the row's baseline or depth.

Make these marks for all global-zero, positive-local tuples of
label \(v\), merging repeated marks at the same tree vertex.
At depth \(\ell\), the number of marks across both trees is at most
the number of used length-\(\lambda_\ell\) local items, hence at
most \(p_\ell\).  All marks are at positive depths.  The retained
depths in Lemma~\ref{lem:competing-trees} ensure that every leaf has
the ancestor used above.

If both counts in Equation~\eqref{eq:sound-side-choice} were less
than \(d-6K\), Lemma~\ref{lem:coverage} would force every chosen
leaf point to be covered by some global-zero, positive-local tuple.
The marks would then hit every root-to-leaf path in both trees,
contrary to the combined-quota obstruction of
Lemma~\ref{lem:competing-trees}.
\end{proof}

\subsection{Job positions and displaced rows}

The forced side at each vertex does not yet ensure that every graph
edge is covered. We must connect the plus allocation to the incident
jobs. A key may use an earlier row, so this connection requires an
explicit count of displaced rows rather than an assumption of
position-by-position matching.

For each label \(v\), exclude a position \(r\in\mathcal R_v\)
if its closed job cell \([b_r-\Delta,b_r]\) meets a coverage interval
from a good \(\mathrm M\) tuple of that label with global length zero
and positive local length.  This definition uses closed cells;
the tested coverage intervals remain half-open.
Such a tuple's interval has length at most
\(\lambda_1+\Delta<g-\Delta\), by
Lemma~\ref{lem:geometry}.  Distinct job cells have gap at least
\(g-\Delta\), so the interval meets at most one of them.
There are at most \(P\) positive local items of label \(v\).
Thus at most \(P\) positions are excluded in all of
\(\mathcal R_v\), and each incident edge retains at least \(R-P\)
nonexcluded positions.

For \(r\in\mathcal R_v\), let \(L_{v,r}\) be the number of good
\(\mathrm G\) tuples whose key is at \(r\) and whose row is a long
row of that same position.

\begin{lemma}[Long matches at a nonexcluded job]\label{lem:job-matches}
Every nonexcluded position \(r\in\mathcal R_v\) satisfies
\begin{equation}\label{eq:sound-long-matches}
 L_{v,r}\ge d-x_v^+-9K.
\end{equation}
\end{lemma}

\begin{proof}
Test at the interior point \(a=b_r-\Delta/2\).
Since position gaps are at least \(g>\Delta\), the job baselines
at most \(a\) are exactly those at positions strictly earlier than
\(r\).  Let \(N'\) count those positions, and let \(\xi\) count
the good tuples whose key is at \(r\) and whose row is from a
strictly earlier position.

The good keys at earlier positions consume at least \(dN'-3K\)
distinct rows from the earlier prefix, by the same order argument
as in Lemma~\ref{lem:coverage}.  The \(\xi\) current-key matches
consume an additional \(\xi\) rows from that prefix: their keys
are distinct from the earlier keys, and a packing cannot use one
row in two tuples.  Among its \(2dN'\) rows, at most
\(dN'+3K-\xi\) can therefore appear in good \(\mathrm M\) tuples.
Including the original tree-row prefix gives the strengthened bound
\begin{equation}\label{eq:sound-displaced-prefix}
 A_v(a)\le B_v(a)+3K-\xi.
\end{equation}

Let \(u\) count the short rows of the current position \(r\) used
in good \(\mathrm M\) tuples.  Every interval covering \(a\)
is accounted for as follows:
\begin{itemize}
 \item Tuples with global length one contribute at most \(x_v^+\)
 in total, since minus tuples cannot reach this plus-side point.
 \item A global-zero tuple with positive local length contributes
 nothing, since the position was not excluded.
 \item With both lengths zero, a tree row or a long job row has
 an empty interval.  A short row covers only its own job interval,
 so it can cover \(a\) only at the current position.
\end{itemize}
It follows that
\begin{equation}\label{eq:sound-job-coverage-upper}
 I_v(a)\le x_v^++u.
\end{equation}
A current short row with global length one may be counted in
both \(x_v^+\) and \(u\); this only enlarges the upper bound.

Combine Equations~\eqref{eq:sound-prefix-decomposition},
\eqref{eq:sound-completion-prefix},
\eqref{eq:sound-displaced-prefix}, and
\eqref{eq:sound-job-coverage-upper} to obtain
\[
 B_v(a)+d-3K
 \le B_v(a)+3K-\xi+x_v^++u.
\]
Equivalently,
\begin{equation}\label{eq:sound-xi-correction}
 u-\xi\ge d-x_v^+-6K.
\end{equation}
Here the same \(3K\) loss controls the entire earlier-key prefix;
it is not incurred separately at its positions.

At least \(d-3K\) keys at \(r\) are in good tuples.
Exactly \(\xi\) of the good matches use earlier rows, and the
row-order restriction forces every other good match to use a row
at \(r\).  At most \(d-u\) of the short rows there are available
to these matches, because \(u\) have already been used in good
\(\mathrm M\) tuples.  Any other use of a short row only lowers
this availability.  Hence
\[
 L_{v,r}\ge d-3K-\xi-(d-u)
          =u-\xi-3K
          \ge d-x_v^+-9K,
\]
as required.
\end{proof}

\subsection{Edge capacity and the threshold cover}

\begin{lemma}[The threshold set covers every edge]\label{lem:edge-cover}
The set
\[
 C=\{v:x_v^+>d/8\}
\]
is a vertex cover.
\end{lemma}

\begin{proof}
Fix an edge index \(e\), with endpoints \(v_1,v_2\).
There are exactly \(2dR(e-1)\)
\(\Waux\) items belonging to earlier edges, counting both permits
and nonpermits.  This equals the total number of earlier-edge
keys over all labels, since each edge has two endpoints.
At least \(2dR(e-1)-3nK\) of those keys are in good tuples.
By the edge-order conclusion of Lemma~\ref{lem:good-feasibility},
each consumes a distinct \(\Waux\) item from the earlier-edge
inventory.  At most \(3nK\) earlier-edge items therefore remain
for all other uses, including good matches of keys on edge \(e\).
This is one global loss bound, independent of the number of
earlier edges.  Uses in other tuples only reduce the remaining
supply.

In an exact-position long match to a key on edge \(e\), the
\(\mathrm G\) inequality becomes
\[
 \beta_f-\beta_e+\Delta\ind\{\text{nonpermit}\}\le0,
\]
where \(f\) is the edge of its \(\Waux\) item.
Lemma~\ref{lem:good-feasibility} gives \(f\le e\);
if \(f=e\), the item must be a permit.  Thus every match counted
by \(L_{v,(e,j)}\) uses either one of the \(dR\) edge-\(e\)
permits or an earlier-edge item.  Summing over both endpoints
and all positions gives
\begin{equation}\label{eq:sound-edge-capacity}
 \sum_{v\in\{v_1,v_2\}}\sum_{j=1}^R L_{v,(e,j)}
 \le dR+3nK.
\end{equation}

Suppose that both endpoints were outside \(C\), so that their
plus counts were at most \(d/8\).
By \(d\ge72K\) and Lemma~\ref{lem:job-matches}, each nonexcluded
position at either endpoint would have
\[
 L_{v,(e,j)}\ge d-\frac d8-9K\ge\frac{3d}{4}.
\]
There are at least \(R-P\) such positions at each endpoint.
Their total demand would therefore be at least
\(2(R-P)(3d/4)\).  Its excess over the upper bound in
Equation~\eqref{eq:sound-edge-capacity} is strictly positive:
\begin{align*}
 2(R-P)\frac{3d}{4}-(dR+3nK)
 &=\frac d2(R-3P)-3nK\\
 &=\frac{97d}{2}P+(50d-3)nK+50d>0,
\end{align*}
using \(R=100(P+nK+1)\) and \(d\ge1\).
This contradiction proves that every edge meets \(C\).
\end{proof}

\begin{lemma}[Cardinality of the extracted cover]\label{lem:cover-size}
The set \(C\) in Lemma~\ref{lem:edge-cover} satisfies
\[
 |C|\le \frac{dk}{d-6K}+\frac{48Kn}{d}
      \le k+\frac{60Kn}{d}
      \le k+\rho n.
\]
\end{lemma}

\begin{proof}
Let
\[
 V_+=\{v:x_v^+\ge d-6K\}.
\]
Since \(d\ge72K\), we have \(d-6K>d/8\), so \(V_+\subseteq C\).
The plus-item budget in Equation~\eqref{eq:sound-global-budgets}
gives
\begin{equation}\label{eq:sound-large-plus-count}
 |V_+|\le\frac{dk}{d-6K}.
\end{equation}

For each vertex in \(V_+\), charge its plus count; for every
other vertex, charge its minus count.
Lemma~\ref{lem:side-choice} makes each of these charges at least
\(d-6K\).  Exactly one side is charged at each vertex, even if
both sides have large counts.  The total charges are therefore
at least \(n(d-6K)\), while the total of all counts is at most
\(dn\).  Every plus count outside \(V_+\) is uncharged, so
\begin{equation}\label{eq:sound-residual-plus}
 \sum_{v\notin V_+}x_v^+\le dn-n(d-6K)=6Kn.
\end{equation}
Each vertex of \(C\setminus V_+\) contributes more than \(d/8\)
to this sum.  The resulting weak bound, valid also when that
set is empty, is
\[
 |C\setminus V_+|\le\frac{48Kn}{d}.
\]
Together with Equation~\eqref{eq:sound-large-plus-count}, this
proves the first asserted inequality.  For the second, use
\(k\le n\) and \(d\ge12K\) to bound
\[
 \frac{dk}{d-6K}-k
 =\frac{6Kk}{d-6K}\le\frac{12Kn}{d}.
\]
The final inequality follows from \(d\ge120K/\rho\).
\end{proof}

\begin{proof}[Proof of Proposition~\ref{prop:soundness}]
Starting from the given packing, Lemma~\ref{lem:edge-cover}
produces the cover \(C\), and Lemma~\ref{lem:cover-size}
bounds its cardinality by \(k+\rho n\).
All deductions used only the stipulated item inventories,
the good-tuple properties, and the prescribed parameters,
so they apply to every packing into at most \(B+c\) bins.
\end{proof}

The polynomial construction in Section~\ref{sec:instance}, integral
completeness in Proposition~\ref{prop:completeness}, and the preceding
soundness proof together establish Theorem~\ref{thm:reduction}.
